\chapter{Transformaciones lineales}
\label{cap:transformaciones}

Una transformación lineal conserva las operaciones del espacio y hace posible comparar espacios distintos mediante una misma estructura. Núcleo e imagen describen su comportamiento algebraico \cite{lang1987,hoffmankunze1971,axler2024}.

\section{Definición y ejemplos}
Linealidad, ejemplos canónicos, proyecciones y transformaciones definidas por fórmulas.

\section{Determinación de una transformación por una base}
Existencia y unicidad de la transformación que asigna imágenes a los vectores de una base.

\section{Núcleo e imagen}
Subespacios asociados y cálculo de sus generadores.

\section{Rango y nulidad}
Dimensiones de la imagen y del núcleo.

\section{Teorema rango--nulidad}
Enunciado, demostración y aplicaciones a la dimensión del dominio.

\section{Inyectividad y sobreyectividad}
Caracterizaciones de inyectividad y sobreyectividad mediante núcleo e imagen.
